We are running this labeling pipeline continuously and have now
collected enough data to evaluate how well it is performing and
study its potential impact on visual recognition.

%We have constructed an initial version of LSUN with
%millions of labeled images in each scene category.
%We experiment with popular convolutional networks using our dataset, and obtain a
%significant performance gain with the same model but trained using our
%bigger training set.

% Given the our pipeline, we want to know the quantity and quality of
% the
% resulting dataset. In addition, we investigate how effectively our
% pipeline is and how the current computer vision models work on our
% dataset.

%%%%%%%%%%%%%%%%%%%%%%%%%%%%%%%%%%%%%%%%

%\vspace{-2mm}
%\subsection{Dataset Statistics}
%\vspace*{-0.05in}

\vspace*{2mm}
\noindent{\bf Dataset Statistics:}
So far, we have collected 59 million images for 20 object categories
selected from PASCAL VOC 2012 and 10 million images for 10 scene
categories selected from the SUN database.
Figure~\ref{fig:numbers} shows the
statistics of object categories. The number of images used in ImageNet
classification challenge is about 1.4 million, which is used to train
recent convolutional networks. Most of our object
categories have more images than the whole challenge. Even if we only
consider the basic level categories in ImageNet such as putting all
the dogs together, our dataset is still much denser.

%% For most of the categories, of which we have far
%% fewer, our dataset is more than 10 times denser (e.g., XXX vs. XXX for
%% bedroom). Similarly, our object categories are approximately 1000 times
%% fewer and denser than ImageNet (e.g., density is XXX vs. XXX for chair).

\begin{figure}[t]
  \centering
    \includegraphics[width=\textwidth]{img/image_numbers.pdf}
   \caption{Number of images in our object categories. Compared to
    ImageNet dataset, we have more images in each category, even
    comparing only basic level categories.}
\label{fig:numbers}
        \end{figure}


%%%%%%%%%%%%%%%%%%%%%%%%%%%%%%%%%%%%%%%%

%\vspace{-2mm}
%\subsubsection{Label Precision}
%\vspace*{-0.05in}

\vspace*{2mm}
\noindent{\bf Label Precision:}
We have run a series of experiments 
to test the labeling precision of this dataset.
We sampled 2000 images from 
11 object categories and label them with 
completely manual pipeline using trained experts (AMT labeling
is not reliable enough for this test).  The final
precision of those categories is shown in
Figure~\ref{fig:pr_precision}a. We observe that normally, the precision
is around 90\%, but it varies across different categories. The major
mistakes in our labels are caused by toys, computer rendered images,
photo collages, and human edited photos.  The inclusion of
computer-rendered images is the most serious for the person
category. Although we have tests for classifier confidence, there are
lots of confusing cases for human labelers on these images,
which subsequently confuse the statistical tests.

\begin{figure}[hbtp]
  \centering
  \begin{subfigure}[b]{0.48\textwidth}
    \includegraphics[width=\textwidth]{img/precision.pdf}    
    \caption{Precision of our labels}
    \label{fig:pr_precision}
  \end{subfigure}
    \begin{subfigure}[b]{0.48\textwidth}
    \includegraphics[width=\textwidth]{img/reduction.pdf}
    \caption{Amplification of human labels}
    \label{fig:pr_reduction}
  \end{subfigure}
\caption{Labeling precision and reduced labeling efforts of our
  pipeline. (a) shows the precision of labels from 11 object
  categories. The precision is normally around 90\%. (b) shows how
  much labeling effort is saved with our
  pipeline.  On average the ratio between all images 
  and the number of human-provided labels is around 40:1.}
\label{fig:precision_reduction}
\end{figure}


%%%%%%%%%%%%%%%%%%%%%%%%%%%%%%%%%%%%%%%%

%\vspace{-2mm}
%\subsubsection{Effort Amplification}
%\vspace*{-0.05in}

\vspace*{2mm}
\noindent{\bf Effort Amplification:}
We next checked how much our pipeline
amplifies human labeling effort -- i.e., the ratio between the number
of images put into a positive or negative set versus the number of images
labeled by people. The result is
shown in Figure~\ref{fig:pr_reduction}. It shows that our pipeline
amplifies human efforts by 40 times on average. In the case of the train
category, people labeled less than 1/60th of the images.


%%%%%%%%%%%%%%%%%%%%%%%%%%%%%%%%%%%%%%%%

%\vspace{-2mm}
%\subsection{Model Performance}
%\vspace*{-0.05in}

\vspace*{2mm}
\noindent{\bf Impact on Model Performance:}
We next investigate whether the new dataset can
help current models achieve better classification performance.
As a first test, we
use the standard AlexNet \cite{AlexNet} trained on PLACES and compare the model fine-tuned
by the PLACES data \cite{zhou2014learning} with the same model fine-tuned
by both PLACES and LSUN, in both cases fine-tuning on only the 10
categories in LSUN.
As a simple measure to balance the dataset,
we only take at most 200 thousand images from each LSUN category.
% \footnote{We didn't ultilize all data that may enable better
% results, due to limited computation resource we have.}. 
Then we test the classification results on
this 10 categories in PLACES testing set. The error percentage comparison is shown in Figure
~\ref{fig:comp_places_error}. The overall detection error percentage
is 28.6\% with only PLACES and 22.2\% with both PLACES and LSUN, 
which yields a 22.37\% improvement on testing error. Although the
testing is unfavorable to LSUN due to potential dataset bias issues
\cite{datasetBias}, we
can see that the additional of more data can improve the classification
accuracy on PLACES testing set for 8 of the 10 categories.

\begin{figure}[hbtp]
\centering
\includegraphics[width=\textwidth]{img/places_error}
\caption{Comparison of classification error by training AlexNet with PLACES
vs. LSUN and PLACES.  Lower errors are better.}
\label{fig:comp_places_error}
\end{figure}

We then use the pre-trained model to fine-tune on
PASCAL VOC 2012 classification training images with Hinge loss and
evaluate the learned representation. We use two pre-trained models for
comparison. One is AlexNet trained by our object data from
scratch. The other is VGG trained by our object data but initialized
by ImageNet. We sample 300 thousand images
from each of the 20 object categories and use them as training data
from LSUN. To get classification score for
testing images, we sample 10
crops from each testing image at center and four corners with mirroring
and take the average. The results are evaluated on the validation
set for each ablation study. As a comparison, we go through the
same procedure with model pre-trained on ImageNet. The results are
shown in Table~\ref{tab:pascal_val}. The table shows that the model
pre-trained on our data performs significantly better. It indicates
that it is better to have more images in relevant categories than more
categories.

Together with the results on scene categories, it is
interesting to observe that although AlexNet is small compared to
other image classification models developed in recent years, even it
can benefit from more training data.

\newcommand\ver[1]{\rotatebox[origin=c]{90}{#1}}
\begin{table}[hbtp]
%\vspace{-2mm}
\begingroup
\setlength{\tabcolsep}{1.5pt}
\scriptsize
\begin{center}
\begin{tabular}{c|c||c|c|c|c|c|c|c|c|c|c|c|c|c|c|c|c|c|c|c|c||c}
Model & Pre-train  & \ver{aero} & \ver{bike} & \ver{bird} & \ver{boat} & \ver{bottle} &
   \ver{bus} & \ver{car} & \ver{cat} & \ver{chair} & \ver{cow} &
     \ver{table} & \ver{dog} & \ver{horse} & \ver{mbike}
   & \ver{person} &
     \ver{plant} & \ver{sheep} & \ver{sofa} & \ver{train} & \ver{tv} &
       \ver{\,mAP} \\ \hline
AlexNet & ImageNet & 0.96 & 0.86 & 0.88 & 0.85 & 0.49 & 0.94 & 0.75 & 0.89 & 0.62 & 0.86 &
   0.65 & 0.87 & 0.92 & 0.9 & 0.85 & 0.53 & 0.91 & 0.68 & 0.92 &
   0.77 & 0.8 \\ 
AlexNet & LSUN &  0.98 & 0.93 & 0.94 & 0.9 & 0.64 & 0.95 & 0.78 & 0.97 & 0.74 & 0.96
   & 0.72 & 0.96 & 0.98 & 0.94 & 0.85 & 0.59 & 0.96 & 0.76 & 0.97 &
   0.82 & 0.87 \\ 
VGG & ImageNet &   0.95 & 0.85 & 0.92 & 0.86 & 0.65 & 0.92 & 0.8 & 0.93 & 0.68 & 0.77
   & 0.66 & 0.9 & 0.85 & 0.86 & 0.94 & 0.56 & 0.89 & 0.6 & 0.94 &
   0.85 & 0.82 \\ 
VGG & LSUN & 0.98 & 0.92 & 0.94 & 0.9 & 0.69 & 0.96 & 0.82 & 0.96 & 0.75 & 0.97
   & 0.78 & 0.93 & 0.97 & 0.92 & 0.95 & 0.67 & 0.96 & 0.73 & 0.97 &
   0.84 & 0.88 \\ \hline
\end{tabular}
\end{center}
\endgroup
\caption{Comparison of pre-training on ImageNet and LSUN. We pre-train
   AlexNet and VGGnet with ImageNet and our data to compare the
   learned representation. We fine tune the networks after removing
   the last layer on PASCAL VOC 2012 images and compare the
   results. The table shows that pre-training with more images in
   related categories is better than that with more categories.}
\label{tab:pascal_val}
\end{table}



%\vspace{-2mm}
%\subsection{Learned Image Representation}
%\vspace*{-0.05in}

\vspace*{2mm}
\noindent{\bf Learned Image Representation:}
As a final test, we study the image representation learned by our object
categories and compare it to ImageNet to understand the tradeoff
between number of images and categories. We compare AlexNet trained by
ImageNet and our training data used in the last section. Then we check the
response in the first layer. As shown in
Figure~\ref{fig:filter}, we find that the
filters in the first level present cleaner patterns than those learned
on ImageNet data.

\begin{figure}[hbtp]
\centering
\includegraphics[width=\textwidth]{img/filter.pdf}
\caption{Comparison of learned filters. We train AlexNet with our
object images and compare the first level filter with ImageNet. We
observe that the filter pattern learned with our data is cleaner,
while the filters from ImageNet are noisier, as highlighted in (b).}
\label{fig:filter}
\end{figure}


%\begin{figure}[hbtp]
%  \centering
%  \begin{subfigure}[b]{\textwidth}
%    \includegraphics[width=\textwidth]{img/places_size}
%    
%    \caption{Comparison of image number for each category. The higher
%      the better.}
%    \label{fig:comp_places_size}
%  \end{subfigure}
%    \begin{subfigure}[b]{\textwidth}
%    \includegraphics[width=\textwidth]{img/places_error}
%    
%    \caption{Comparison of classification error by training
%      AlexNet. The
%      lower the better.}
%    \label{fig:comp_places_error}
%  \end{subfigure}
%\caption{Comparison of LSUN and PLACES. (a) For most of the
%  categories, we
%are more than 10 times bigger than PLACES. Please note that the scale
%of y axis is log. (b) We train AlexNet with PLACES and putting PLACES
%and LSUN together. With addition of LSUN data, the trained model has
%better
%performance evaluated by PLACES test set.}
%\label{fig:comp_places}
%\end{figure}



%% \begin{table}[t]
%% \small
%%   \centering
%%   \begin{tabular}{ c | c | c | c | c | c }
%%     \hline
%%     Iteration & Method & Positive images & Precision & Positive labels
%%     & Label Ratio \\ \hline
%%     0 & Clustering & 941,981 & 96.9\% & 9,515 & 1\%\\ \hline
%%     1 & ConvNet Feature + SVM & 1,918,913 & 96.4\% & 41,785 & 2.1\%
%%     \\ \hline
%%     2 & ConvNet Feature + SVM & 2,011,332 & 96.2\% & 88,259 & 4.4\%
%%     \\ \hline
%%     3 & ConvNet Fine Tuning & 2,140,763 & 96.1\% & 88,259 & 4.1\%
%%     \\ \hline
%%     4 & ConvNet Fine Tuning & 2,215,244 & 95.9\% & 147,453 & 6.6\%
%%     \\ \hline
%%   \end{tabular}
  
%%   \vspace{-2mm}
%%   \caption{Algorithm statistics of kitchen. The third column shows
%%     the cumulative number of positive images in each iteration. The
%%     fifth column shows the cumulative number of positive labels
%%     provided by human. The last column shows the ratio of human
%%     labeled images compared to all the positive images we
%%     obtained. See the main text for more details.}
%%   \vspace{-3mm}
%%   \label{tab:alg_stat}
%% \end{table}

